\documentclass[11pt,a4paper]{ctexart}

\usepackage{amsmath,amssymb,mathtools}
\usepackage{geometry}
\usepackage{booktabs}
\usepackage{hyperref}
\usepackage{xcolor}
\usepackage{enumitem}

\geometry{margin=2.35cm}
\hypersetup{colorlinks=true, linkcolor=blue!60!black, urlcolor=blue!60!black}
\setlist[itemize]{leftmargin=2em}
\setlist[enumerate]{leftmargin=2em}

\title{四种生成范式的 Tensor Flow 视角：自回归、扩散、流匹配与一致性模型}
\author{}
\date{\today}

\begin{document}
\maketitle
\tableofcontents
\newpage

\section{为什么用 Tensor Flow 视角}

这份笔记不从抽象概率论开始，而从模型内部的张量流动开始。我们关心四个问题：
\begin{enumerate}
  \item 输入张量是什么 shape？
  \item backbone 输出什么 latent / hidden state？
  \item head 把 latent 解码成什么？
  \item loss 在哪些位置、哪些维度上计算？
\end{enumerate}

统一记号如下：
\[
  B=\text{batch size},\quad
  N=\text{token 数或序列长度},\quad
  D=\text{hidden dimension},\quad
  V=\text{vocabulary size}.
\]
在 LLM/VLA 中常见张量为：
\[
  H\in\mathbb{R}^{B\times N\times D}.
\]
其中 \(H[:,i,:]\) 是第 \(i\) 个 token 位置的 hidden state。

\section{自回归：hidden state 经过 LM head 变成下一个 token}

\subsection{核心张量流}

自回归模型的输入是一串 token：
\[
  X=[x_1,x_2,\ldots,x_N]\in\mathbb{N}^{B\times N}.
\]
经过 embedding 和 Transformer：
\[
  X
  \xrightarrow{\mathrm{Embed}}
  E\in\mathbb{R}^{B\times N\times D}
  \xrightarrow{\mathrm{Transformer}}
  H\in\mathbb{R}^{B\times N\times D}.
\]
LM head 是一个线性层：
\[
  W_{\mathrm{lm}}\in\mathbb{R}^{D\times V}.
\]
因此：
\[
  \mathrm{logits}=H W_{\mathrm{lm}}
  \in\mathbb{R}^{B\times N\times V}.
\]

这句话非常关键：\textbf{每一个 token 位置都会输出一个完整词表上的 logits}。

\subsection{训练时：并行预测所有 next token}

训练时使用 teacher forcing。模型输入是：
\[
  [x_1,x_2,\ldots,x_{N-1}],
\]
目标是：
\[
  [x_2,x_3,\ldots,x_N].
\]
所以第 \(i\) 个位置的 hidden state 用来预测第 \(i+1\) 个 token：
\[
  H[:,i,:]\rightarrow \mathrm{logits}[:,i,:]\rightarrow x_{i+1}.
\]
交叉熵损失为：
\[
  \mathcal{L}_{\mathrm{AR}}
  =
  \mathrm{CE}\left(
    \mathrm{logits}[:,0:N-1,:],
    X[:,1:N]
  \right).
\]

\subsection{推理时：只用最后一个位置}

推理时，当前序列长度为 \(N\)。Transformer 输出：
\[
  H\in\mathbb{R}^{B\times N\times D}.
\]
生成下一个 token 只取最后一个位置：
\[
  h_{\mathrm{last}}=H[:,-1,:]\in\mathbb{R}^{B\times D}.
\]
经过 LM head：
\[
  \mathrm{logits}_{\mathrm{next}}
  =
  h_{\mathrm{last}}W_{\mathrm{lm}}
  \in\mathbb{R}^{B\times V}.
\]
然后从 \(V\) 个 token 中选择一个：
\[
  x_{N+1}\sim\mathrm{Softmax}(\mathrm{logits}_{\mathrm{next}}).
\]

\subsection{OpenVLA 的 action token 例子}

OpenVLA 原版没有单独的 continuous action head。它把连续动作离散化为 action tokens，然后仍然用 LM head 预测 token。

假设动作是 7 维：
\[
  a=[a^1,a^2,\ldots,a^7].
\]
每一维动作先被量化成一个 action token：
\[
  a^j \rightarrow y_j,\qquad y_j\in\{\text{action token ids}\}.
\]
推理时不是一次输出 7 维连续向量，而是自回归生成 7 个 token：
\[
  y_1,y_2,\ldots,y_7.
\]
每一步都是：
\[
  H[:,-1,:]\in\mathbb{R}^{B\times D}
  \rightarrow
  \mathrm{LMHead}
  \rightarrow
  \mathbb{R}^{B\times V}
  \rightarrow
  y_j.
\]
最后 ActionTokenizer 把 7 个 action tokens 反量化为连续动作：
\[
  [y_1,\ldots,y_7]\rightarrow [a^1,\ldots,a^7].
\]

\subsection{真正的 continuous action head}

如果不用 action token，而是直接输出连续动作，则会接一个 action head：
\[
  h_{\mathrm{last}}\in\mathbb{R}^{B\times D}
  \xrightarrow{\mathrm{MLP}}
  \hat{a}\in\mathbb{R}^{B\times A},
\]
其中 \(A\) 是动作维度，例如 \(A=7\)。这时 loss 可以是：
\[
  \mathcal{L}_{\mathrm{act}}
  =
  \|\hat{a}-a\|_1
  \quad\text{或}\quad
  \|\hat{a}-a\|_2^2.
\]

\section{扩散：把 clean action/video 张量加噪，再预测噪声}

\subsection{生成对象的 shape}

扩散模型处理的是连续张量。比如 WAM 中要生成一个 action chunk：
\[
  x_0=a_{t:t+H-1}\in\mathbb{R}^{B\times H\times A}.
\]
如果生成未来视频 latent 和动作的联合对象，可以写成：
\[
  x_0=[z^o,a]\in\mathbb{R}^{B\times L\times C}.
\]
这里 \(L\) 是 token/patch/chunk 长度，\(C\) 是每个 token 的通道维。

\subsection{训练时的张量流}

训练时先采样时间步 \(t\)，再采样噪声：
\[
  \epsilon\sim\mathcal{N}(0,I),
  \qquad
  \epsilon\in\mathbb{R}^{B\times L\times C}.
\]
把 clean sample 加噪：
\[
  x_t
  =
  \alpha_t x_0+\sigma_t\epsilon.
\]
模型输入通常是：
\[
  (x_t,t,c),
\]
其中 \(c\) 是条件，例如图像观测、语言指令、proprioception。

denoising network 输出：
\[
  \hat{\epsilon}
  =
  \epsilon_\theta(x_t,t,c)
  \in\mathbb{R}^{B\times L\times C}.
\]
注意：\textbf{输出 shape 与噪声 shape 完全一样}。

loss 是逐元素回归：
\[
  \mathcal{L}_{\epsilon}
  =
  \|\hat{\epsilon}-\epsilon\|_2^2.
\]

\subsection{推理时的张量流}

推理从纯噪声开始：
\[
  x_T\sim\mathcal{N}(0,I),
  \qquad
  x_T\in\mathbb{R}^{B\times L\times C}.
\]
每一步：
\[
  x_t
  \xrightarrow{\epsilon_\theta(x_t,t,c)}
  \hat{\epsilon}
  \xrightarrow{\mathrm{scheduler}}
  x_{t-1}.
\]
重复多步：
\[
  x_T\rightarrow x_{T-1}\rightarrow\cdots\rightarrow x_0.
\]

\subsection{和 action head 的区别}

扩散模型的输出 head 不是 vocabulary logits，而是一个和 \(x_t\) 同 shape 的连续张量：
\[
  \mathbb{R}^{B\times L\times C}\rightarrow
  \mathbb{R}^{B\times L\times C}.
\]
它不是在词表里选 token，而是在每个连续维度上预测噪声、velocity 或 clean sample。

\section{流匹配：同样输入中间张量，但预测 velocity}

\subsection{训练样本的构造}

Flow Matching 也处理连续张量。设：
\[
  x_0\sim p_{\mathrm{noise}},
  \qquad
  x_1\sim p_{\mathrm{data}},
\]
且二者 shape 相同：
\[
  x_0,x_1\in\mathbb{R}^{B\times L\times C}.
\]
构造中间状态：
\[
  x_t=(1-t)x_0+t x_1.
\]
真实速度为：
\[
  u_t=x_1-x_0.
\]

\subsection{模型输入输出}

模型输入：
\[
  (x_t,t,c).
\]
模型输出 velocity：
\[
  \hat{v}
  =
  v_\theta(x_t,t,c)
  \in\mathbb{R}^{B\times L\times C}.
\]
loss：
\[
  \mathcal{L}_{\mathrm{FM}}
  =
  \|\hat{v}-(x_1-x_0)\|_2^2.
\]

所以从 tensor flow 上看，Flow Matching 和 diffusion 很像：
\[
  \text{input: noisy/intermediate tensor}
  \rightarrow
  \text{network}
  \rightarrow
  \text{same-shape continuous tensor}.
\]
区别是 target 不同：
\[
  \text{diffusion target}=\epsilon,
  \qquad
  \text{flow matching target}=u_t.
\]

\subsection{推理时的张量流}

生成从噪声开始：
\[
  x(0)\in\mathbb{R}^{B\times L\times C}.
\]
每个 ODE step：
\[
  x(t)
  \xrightarrow{v_\theta(x(t),t,c)}
  \hat{v}(t)
  \xrightarrow{\mathrm{ODE\ solver}}
  x(t+\Delta t).
\]
最终得到：
\[
  x(1)\approx x_{\mathrm{data}}.
\]

\section{一致性模型：从任意噪声时间点直接映射到 clean sample}

\subsection{模型输入输出}

Consistency Model 的输入仍是某个时间点的 noisy/intermediate tensor：
\[
  x_t\in\mathbb{R}^{B\times L\times C}.
\]
但模型输出不是噪声，也不是速度，而是 clean sample 的估计：
\[
  \hat{x}_0=f_\theta(x_t,t,c)
  \in\mathbb{R}^{B\times L\times C}.
\]

\subsection{一致性约束}

取同一条轨道上的两个时间点 \(s<t\)：
\[
  x_s,\ x_t.
\]
模型应满足：
\[
  f_\theta(x_t,t,c)
  \approx
  f_{\theta^-}(x_s,s,c).
\]
loss：
\[
  \mathcal{L}_{\mathrm{CM}}
  =
  d\left(
    f_\theta(x_t,t,c),
    f_{\theta^-}(x_s,s,c)
  \right).
\]
这里 \(f_{\theta^-}\) 常是 EMA teacher 或 stop-gradient target。

\subsection{推理}

一步生成：
\[
  x_T\sim\mathcal{N}(0,I),
  \qquad
  \hat{x}_0=f_\theta(x_T,T,c).
\]
少步生成则是：
\[
  x_T\rightarrow \hat{x}_{t_1}\rightarrow \hat{x}_{t_2}\rightarrow \hat{x}_0.
\]

\section{四种范式的 Head 对比}

\begin{center}
\begin{tabular}{llll}
\toprule
范式 & backbone 输出 & head 输出 & 输出 shape \\
\midrule
自回归 & hidden states \(H\) & vocab logits & \(B\times N\times V\) \\
连续动作 AR & last hidden state & action vector & \(B\times A\) 或 \(B\times H\times A\) \\
扩散 & noisy latent features & noise / velocity / clean tensor & \(B\times L\times C\) \\
流匹配 & intermediate latent features & velocity tensor & \(B\times L\times C\) \\
一致性 & noisy/intermediate latent features & clean tensor estimate & \(B\times L\times C\) \\
\bottomrule
\end{tabular}
\end{center}

\section{以 WAM action chunk 为例的统一对比}

假设要生成一段机器人动作：
\[
  a_{t:t+H-1}\in\mathbb{R}^{B\times H\times A}.
\]

\subsection{自回归 token 版本}

先把每个动作维度离散化：
\[
  a_{h,j}\rightarrow y_{h,j}\in\{1,\ldots,V\}.
\]
模型生成 token 序列：
\[
  y_1,y_2,\ldots,y_{H\cdot A}.
\]
每一步：
\[
  H[:,-1,:]\rightarrow \mathrm{LMHead}\rightarrow \mathbb{R}^{B\times V}.
\]
最后 detokenize：
\[
  y_{1:H\cdot A}\rightarrow a_{t:t+H-1}.
\]

\subsection{连续 action head 版本}

Transformer 得到最后 hidden state 或 pooled representation：
\[
  h\in\mathbb{R}^{B\times D}.
\]
action head 直接输出：
\[
  \hat{a}
  =
  \mathrm{MLP}(h)
  \in\mathbb{R}^{B\times H\times A}.
\]
这是并行输出整个 action chunk。

\subsection{扩散版本}

从 noisy action chunk 开始：
\[
  a_t^{\mathrm{noisy}}\in\mathbb{R}^{B\times H\times A}.
\]
denoising network 每一步输出同 shape 噪声：
\[
  \hat{\epsilon}\in\mathbb{R}^{B\times H\times A}.
\]
多步去噪得到：
\[
  \hat{a}_{t:t+H-1}.
\]

\subsection{流匹配版本}

从 noise action tensor 到 data action tensor 的中间状态：
\[
  a_\tau=(1-\tau)a_0+\tau a_1.
\]
网络输出 velocity：
\[
  \hat{v}\in\mathbb{R}^{B\times H\times A}.
\]
通过 ODE solver 积分得到 action chunk。

\subsection{一致性版本}

输入任意噪声时间点的 action tensor：
\[
  a_\tau\in\mathbb{R}^{B\times H\times A}.
\]
模型直接输出 clean action chunk：
\[
  \hat{a}=f_\theta(a_\tau,\tau,c)
  \in\mathbb{R}^{B\times H\times A}.
\]

\section{一句话总结}

从 tensor flow 视角看，自回归的核心是
\[
  B\times N\times D
  \rightarrow
  B\times N\times V,
\]
即每个 token 位置通过 LM head 预测下一个 token。扩散、流匹配和一致性模型的核心则是
\[
  B\times L\times C
  \rightarrow
  B\times L\times C,
\]
区别只在 target：扩散预测噪声，流匹配预测速度，一致性模型预测 clean sample 或轨道终点。

\end{document}
